%! Representing a Quantitative Variable with Graphs — Unit 1, Lesson 1.3 — Lesson Plan
% GRADUAL-RELEASE plan (warm-up / guided notes and practice / debrief / close and start homework).
% THERE IS NO GROUP ACTIVITY IN THIS COURSE — the guided notes carry the release: I do (two
% sections) -> we do (guided practice). THERE IS NO INDEPENDENT PRACTICE SET — the homework is
% the individual practice, and the period ends with students starting it. Teacher-facing. Every box is
% `skillbox` (breakable) with a \boxguard ahead of it; this course has no `fixedskillbox`,
% no tier boxes, and no exit ticket.
\documentclass[10pt]{article}
\usepackage{apstats-article}
\usepackage{apstats-boxes}

\newcommand{\UnitNumberName}{Unit 1: Exploring One-Variable Data \quad}
\newcommand{\LessonNumberName}{Lesson 1.3: Representing a Quantitative Variable with Graphs}

\begin{document}
\begin{center}
    {\Large\bfseries \CourseName \\
    {\small\bfseries \UnitNumberName \LessonNumberName}}
\end{center}

\begin{tcolorbox}[colback=sky, colframe=navy, boxrule=0.9pt, arc=2mm,
    left=2mm, right=2mm, top=1.5mm, bottom=1.5mm]
    \textbf{Primary Objective:} Students will classify a quantitative variable as discrete or
    continuous, read dotplots, stem plots, and histograms, and explain how changing the interval
    width of a histogram changes its appearance without changing a single data value.
    \\[2pt]
    \textbf{CED alignment:} Topic 1.5 \quad$\cdot$\quad Big Idea: UNC (Uncertainty)
    \quad$\cdot$\quad Skills 2.A, 2.B \quad$\cdot$\quad LO UNC-1.F, UNC-1.G \quad$\cdot$\quad
    EK UNC-1.F.1, 1.F.2, 1.G.1--1.G.5
    \\[2pt]
    \textbf{Lesson model:} \emph{Gradual release --- I do, we do, you do --- all of it inside
    the guided notes.} There is no group activity. The notes name the vocabulary as they teach
    it and deliver the instruction in two sections; the guided practice is worked together with
    students holding the pen; \textbf{the homework is the individual practice}, and students
    start it in class while you circulate. The whole-class debrief is \emph{spoken} ---
    there is no transfer set and no reflection box on the page.
    \\[2pt]
    \textbf{Note on the ``no sketch'' rule:} every display in this lesson is \textbf{given}.
    Students construct only by filling in the interval table in notes section 2. Do not add a
    ``now draw your own histogram'' step --- the AP exam asks students to \emph{read} and
    \emph{choose} displays, and drawing axes by hand consumes the period without teaching the
    objective.
\end{tcolorbox}

\renewcommand{\arraystretch}{1.4}
\boxguard
\begin{skillbox}[Learning Targets \& Key Understandings]{goldbox}
\begin{tabularx}{\linewidth}{@{}X|X@{}}
\textbf{Learning targets (student language)} & \textbf{Key understandings (the ``why'')} \\[2pt]
\hline
\vspace{-6pt}
\begin{itemize}[leftmargin=1.1em, nosep, after=\vspace{2pt}]
  \item I can say which measurements can land between two whole numbers and which cannot.
  \item I can read a dotplot, a stem plot, and a histogram, and say what each one keeps.
  \item I can explain what grouping data gains and what it throws away.
  \item I can explain how two correct pictures of the same numbers can disagree.
\end{itemize}
&
\vspace{-6pt}
\begin{itemize}[leftmargin=1.1em, nosep, after=\vspace{2pt}]
  \item Discrete variables take a countable set of values; continuous variables take
        infinitely many, with another value always available between any two
        (EK UNC-1.F.1, 1.F.2).
  \item Dotplots and stem plots preserve every individual value; a histogram trades those
        values for interval counts (EK UNC-1.G.1--1.G.3).
  \item \textbf{Target misconception:} that a histogram reports a fact about the data. It
        reports a fact about the data \emph{and the interval width chosen}. The width-5 picture
        shows two clumps and a gap; the width-20 picture of the identical 24 values shows two
        equal bars and no structure at all (EK UNC-1.G.1).
  \item \textbf{Second misconception:} that ``the values are whole numbers'' makes a variable
        discrete. Rounded measurements are still continuous --- ask about the variable, not the
        recorded digits.
  \item \textbf{Third misconception:} the over-correction --- that narrower blocks are always
        better. At width 1 these 24 times give a bar of height 0 or 1 almost everywhere.
\end{itemize}
\end{tabularx}
\end{skillbox}

\boxguard
\begin{skillbox}[{Vocabulary, Concepts, \& Theorems --- taught directly in the guided notes}]{greenbox}
\begin{minipage}[t]{0.48\linewidth}
    \begin{itemize}
        \item \textbf{Discrete variable} --- a countable number of values (EK UNC-1.F.1)
        \item \textbf{Continuous variable} --- infinitely many values; another always lies
              between any two (EK UNC-1.F.2)
        \item \textbf{Dotplot} --- one dot per observation, stacked (EK UNC-1.G.3)
        \item \textbf{Stem plot} --- value split into stem (leading digits) and leaf (last
              digit) (EK UNC-1.G.2)
    \end{itemize}
\end{minipage}
\hfill
\begin{minipage}[t]{0.48\linewidth}
    \begin{itemize}
        \item \textbf{Histogram} --- bar height is the count in an interval; bars touch
              (EK UNC-1.G.1)
        \item \textbf{Interval width} --- the choice that controls what the histogram shows
        \item \textbf{Cumulative graph} --- how many observations are at or below a value
              (EK UNC-1.G.4) --- named in section 2, worked in Lesson 1.6
        \item \textbf{Carried forward} --- categorical vs.\ quantitative (1.1); bar graphs
              and their gaps (1.2)
    \end{itemize}
\end{minipage}
\end{skillbox}

% A tabularx never splits, so guard generously ahead of it.
\boxguard[30]
\begin{skillbox}[Lesson at a Glance --- \MeetingLength]{sky}
\renewcommand{\arraystretch}{1.25}
\begin{tabularx}{\linewidth}{@{}l c X X@{}}
\textbf{Phase} & \textbf{Min} & \textbf{Students} & \textbf{Teacher} \\
\hline
Warm-Up & 5 & Shares, a tally, and four things that might be 7.5 & Debrief item~3 aloud; leave their ``you can't have half a car'' sentence on the board all period \\
Guided Notes \& Practice & 34 & Fill the vocabulary box and notes 1--2; then \emph{The Bus Stop} together, pen in their hand & \textbf{I do} notes 1--2 (20) $\to$ \textbf{we do} the Guided Practice box (14) \\
Debrief & 8 & Pens down, papers up --- answer aloud, nothing to fill in & Put both histogram totals side by side, then take the rounding trap and the narrow-is-better over-correction \\
Close \& Assign & 8 & \textbf{Start the homework in class}, alone and silent & Announce the paper homework, circulate the first minutes for your formative read, preview Lesson~1.4 \\
\end{tabularx}
\end{skillbox}

\boxguard[20]
\begin{skillbox}[Warm-Up --- Activate Prior Knowledge (5 min)]{sky}
\begin{minipage}[t]{0.48\linewidth}
    \textbf{The three items and what each seeds}
    \begin{itemize}
        \item \textbf{1.} Two nights of ticket sales as shares. \textbf{Spirals} to Lesson 1.2
              --- and quietly confirms that today's variable is the other kind.
        \item \textbf{2.} Tally twelve numbers into three blocks. \textbf{Seeds:} the exact move
              notes section~2 asks for, done once with easy numbers first.
        \item \textbf{3.} Which of these could be 7.5? \textbf{Seeds:} discrete vs.\ continuous,
              in language students already own --- section~1 names it fifteen minutes later.
    \end{itemize}
\end{minipage}
\hfill
\begin{minipage}[t]{0.48\linewidth}
    \textbf{Running it}
    \begin{itemize}
        \item \textbf{Debrief item~3 aloud} and write their own phrasing on the board
              (``you can't have half a car''). Leave it up: section~1 attaches
              \emph{discrete} and \emph{continuous} to that sentence rather than to a
              definition you supply cold.
        \item Expect an argument about time: \emph{``a stopwatch only shows hundredths.''}
              That is the rounding misconception arriving early, which is good --- name the
              two positions, do not settle it, and tell them section~1 decides it.
        \item Item~2 is mechanical. If the room is quick, take 90 seconds and move on. It
              exists so that grouping is not the obstacle inside section~2.
    \end{itemize}
\end{minipage}
\end{skillbox}

\boxguard
\begin{skillbox}[Guided Notes \& Practice --- I do, then we do (34 min)]{sky}
\textbf{This is the whole lesson.} There is no group activity, and there is no independent
practice set on this page: \textbf{the homework is the individual practice}, started in class.
\begin{multicols}{2}
\textbf{I do (about 20 min).} Two sections, each in two moves; students fill the blanks as each
idea is named, and the vocabulary box is filled \emph{as you go}, never in advance.

\emph{Section 1} opens on the two kinds of variable, hung on the warm-up sentence still on the
board. Finish that half on the red callout --- a gauge reading 13\,mm has not made rainfall
discrete. About five minutes. Its second half, \emph{The Class Record}, puts the stem plot and
the dotplot of the same 24 finishing times side by side; the point is not the mechanics but that
\emph{two displays that look nothing alike carry identical information}, because both keep every
value. Make them find the two 36s in each. Four minutes.

\textbf{Section 2 is the lesson.} Its first half is the hinge: they fill the interval table
themselves, and then you say the line that makes the histogram cheap --- \textbf{that table is
already the histogram}. Take the bars-touch contrast straight back to Lesson 1.2's gaps, and
name the cumulative graph in one sentence without working it. Five minutes. Its second half,
\emph{The Same 24 Times, Drawn Twice}, is the crux: put both pictures up before computing
anything and ask which one is lying. Then work the \texttt{work} block --- both add to 24 --- so
that ``data went missing'' is off the table before the argument starts. The sentence to land is
\emph{a histogram is a choice, not a fact}; say it, write it, and point at the two pictures
while you do. Close on the counterweight: blocks of one minute give every value back and no
shape at all.

\textbf{We do (about 14 min).} The Guided Practice box, \emph{The Bus Stop} --- a third display
in a fresh context. Students hold the pen; you hold the questions: \emph{``What makes those two
riders look strange --- where they are, or what is next to them?''} \quad \emph{``Draw me the
5--10 block. Is it empty?''} Part (c) is the crux transferred and is the part to protect when
time is short.

\columnbreak
\textbf{The pen must actually change hands.} Guided Practice is the only place today where you
can watch a student decide, so do not narrate it. Put part (a) on them cold, take a hand count on
\emph{discrete or continuous} in part (b) before anyone speaks, and make the room defend both
answers.

Circulate during Guided Practice and demand the reason, not the letter:
\begin{itemize}[leftmargin=1.1em, nosep]
  \item \textbf{Part (a):} ``Longer than 8 minutes --- read me the dots you counted, not the
        number.''
  \item \textbf{Part (b), the rounding trap:} ``The picture only shows whole minutes. Did the
        bus care?''
  \item \textbf{Part (c), the crux:} ``Draw me the 5--10 block. Is it empty? Then what happens
        to the two riders out at 11 and 12?''
\end{itemize}
While circulating, pick the papers to put up in the debrief --- one whose part (c) kept the two
riders but lost the empty stretch, one that missed that the stretch went with them.
\textbf{Your formative read comes twice}: here, and again in the last eight minutes as students
start the homework.
\end{multicols}
\end{skillbox}

\boxguard
\begin{skillbox}[Debrief --- whole class (8 min)]{sky}
Pens down, papers up. \textbf{This debrief is spoken} --- there is no box at the end of the
notes to fill in, so it lives entirely on the board and in the room.
\begin{multicols}{2}
\begin{enumerate}[leftmargin=1.3em, nosep, itemsep=3pt]
  \item \textbf{Put both totals on the board first:} 24 and 24. Nothing was thrown away by
        either picture, so the only remaining explanation is the width. Then state it as
        EK UNC-1.G.1: altering the interval width changes the appearance.
  \item Put up the two Guided Practice papers you picked and name what the correct one
        did. Part (c) regroups in the \emph{other} direction from the notes, so this is where
        you find out whether students learned the rule or learned ``wide is bad.''
  \item Circle the warm-up sentence still on the board and take the rounding trap off GP~(b):
        the question is about the variable, not the digits recorded.
  \item Take the over-correction last: \emph{``How many of the 24 times repeat? So how tall
        would most one-minute bars be?''}
\end{enumerate}
\columnbreak
\textbf{Two things to demand aloud.} First, \textbf{make a student say what a histogram of the
24 times does and does not support} --- seven students between 15 and 19 minutes, and not one of
those seven times recoverable. Second, \textbf{make them name a question that a dotplot answers
and a histogram cannot}, and then one where 200 values make the histogram the better choice.
Without the second half, students leave believing histograms are the inferior display. Hold the
room to both rather than letting students race ahead to the homework --- they will start it in
eight minutes.

\textbf{Connect back to Lesson~1.2 if time allows:} the bar graph's gaps and the histogram's
touching bars are the same fact about the horizontal axis, read twice. Categories are names;
intervals are stretches of a number line.

\textbf{If the debrief runs short}, cut the Lesson~1.2 callback and the rounding trap --- item~1
is the only one that cannot be cut. \textbf{Do not let the debrief eat the last eight minutes}
--- the homework start is not optional padding, it is your second formative read.
\end{multicols}
\end{skillbox}

\boxguard
\begin{skillbox}[AP Practice --- assigned, not scored]{goldbox}
Four AP-style multiple-choice items on a 20-day rainfall stem plot --- a counting item, the
discrete/continuous judgment with the ``values are whole numbers'' distractor, the
regroup-to-a-histogram item, and the crux as a both-are-correct item --- then one multi-part free
response on the same record: classify, group into 10\,mm blocks, state what grouping costs, and
reject a bar graph. \textbf{Parts (a) and (d) are the spiral} --- (a) reaches back to Lesson 1.1
(categorical vs.\ quantitative) and (d) to Lesson 1.2 (bar graphs are for categories). Reusing one
display for both sections is deliberate: the context is already loaded, so the thinking stays on
the choices.

\textbf{Not scored --- the cover's score column reads \textbf{NA} for this page.} The rainfall
record appears nowhere else in the packet, so what comes back in Lesson~1.4 is your evidence
that the idea transferred to a data set students met on their own. Go over it at the start of
Lesson~1.4 and hold the AP standard out loud: \emph{in context or it does not count}. A part (c)
that says ``the table loses information'' without naming \emph{which} information is incomplete;
the answer must point at something concrete, such as the two days that both recorded 13\,mm.
\end{skillbox}

\boxguard
\begin{skillbox}[Homework --- scored, due the first class after two study halls]{goldbox}
Five exercises spanning Topic 1.5 in fresh contexts: \textbf{1} discrete/continuous on four
variables; \textbf{2} the crux restated --- two histograms of 60 test scores, and \emph{neither}
is wrong; \textbf{3} the bars-touch/bars-gap contrast, which is the link back to Lesson 1.2;
\textbf{4} reading a stem plot row and saying why a histogram cannot return those values;
\textbf{5} an AP-style multiple-choice item with a required justification.

\textbf{Start it in the last eight minutes of the period, in class and alone.} \textbf{Scoring:} 10 points --- 2 per
item. \textbf{Item~2 is where the misconception shows}: students who answer ``the 10-point one is
wrong'' have heard ``wide hides things'' and turned it into a rule. The answer is \emph{neither}.
\textbf{Item~5 carries the check:} distractor (E) is that same over-correction in AP clothing.

\textbf{Paper homework is the default for this course} --- assign this page unless you have
decided otherwise. \textbf{DeltaMath override (occasional):} on the days you assign DeltaMath
instead, it replaces this page, you strike through the score cell on the cover, and this page
stays in the packet as extra practice. Never assign both.
\end{skillbox}

\boxguard
\begin{skillbox}[Watch For (while circulating)]{redbox}
\begin{itemize}[leftmargin=1.2em, nosep]
  \item \textbf{``One of these pictures must be wrong''} (notes~2, GP~(c), HW~2 --- the target
        misconception). Probe: \emph{``Add both up. Same 24? Then which one lied?''}
  \item \textbf{Whole numbers read as discrete} (notes~1, GP~(b), HW~1, AP~MC2). Rainfall
        recorded as 13\,mm is still continuous. Probe: \emph{``Could it have rained 13.4 mm?
        Would the gauge have cared?''}
  \item \textbf{Over-correction to ``narrow is always better''} (notes~2, HW~5, AP~MC4):
        the mirror error, and the one the AP exam actually tests. Probe: \emph{``How many of the
        24 times repeat? Draw me the one-minute picture.''}
  \item \textbf{Histogram bars drawn or read with gaps} (notes~2, HW~3): a leftover from Lesson
        1.2. Probe: \emph{``What is on the horizontal axis here --- names, or numbers?''}
  \item \textbf{Values invented back out of a histogram} (notes~2, HW~4, AP~c): students
        report ``a student finished in 37 minutes'' from a 35--39 bar. Probe: \emph{``Point at
        the 37.''}
  \item \textbf{Stem plot misread} as 3 and 4 rather than 34 (notes~1, HW~4). Probe:
        \emph{``Read me the whole number, not the two digits.''}
\end{itemize}
Cold-call prompts: \emph{``Name a variable that is quantitative and discrete.''} \quad
\emph{``What would I have to change to make this histogram show two peaks?''}
\end{skillbox}

\boxguard
\begin{skillbox}[Close \& Assign (8 min)]{goldbox}
\textbf{Students start the homework here, in class, alone and silent --- this is the individual
practice, and the first minutes of it are your best formative read.} Hand the launch first ---
five exercises on paper, scored, due the first class after two study halls --- and point
out that the AP Practice page before it is theirs to work and bring to review. Circulate:
\textbf{item~2 is the column that tells you whether section~2 landed} --- two histograms of the
same 60 test scores, and \emph{neither} is wrong. Sort what you see into three piles as you walk
--- said \emph{neither} and named the width (ready); said \emph{neither} with no reason (ready,
needs the language); named one version wrong (the misconception survived) --- and open
Lesson~1.4 by reteaching section~2 if that third pile ran past a third of the room. Then close
on what changed: \emph{``You walked in thinking a graph shows you the data. You walk out knowing a
histogram shows you the data \emph{and} the choice somebody made about it --- and that you have
to check that choice before you believe the shape.''} \textbf{Preview:} Lesson~1.4,
\emph{Describing the Distribution of a Quantitative Variable} --- now that you can read these
displays, we put words to them: shape, centre, spread, and unusual values, in context, every
time. Note that \emph{cumulative graphs} (EK UNC-1.G.4) are named today but not worked; they
return in Lesson~1.6 alongside percentiles, where they carry their weight.
\end{skillbox}

% ── Teacher notes — one per component, in packet order (four: there is no activity) ──
\begin{teachernote}[Warm-Up]
Item~3 is the seed and item~2 is the drill. Give four minutes total. Expect the time argument on
item~3 --- \emph{``a stopwatch only shows hundredths, so time is countable''} --- and do not
settle it here; write both positions on the board and tell them section~1 of the notes decides
it. That suspense is worth more than a correct answer now, and the red callout in section~1
resolves it cleanly. \textbf{Write their own phrasing on the board verbatim and leave it up all
period}; you will circle it in the debrief. Item~1 is a 60-second spiral to Lesson 1.2 --- do not
let it expand.
\end{teachernote}

\begin{teachernote}[Guided Notes \& Practice]
\textbf{This one document is the lesson --- pace it 20 / 14, then 8 for the debrief, then 8 for
the homework start.} Section~1 is the machinery and must move: five minutes on the two kinds of
variable, four on the displays that keep every value. The temptation is to teach the stem plot
carefully; resist it. The second half of section~1 exists to establish one thing --- both
displays keep every value --- and the two 36s prove it faster than any explanation.

\textbf{The first half of section 2 is the hinge, and the interval table is the reason it
works.} Students fill it in themselves, and then the histogram costs you one sentence:
\emph{that table is already the histogram}. If you draw the histogram first and the table
second, the crux lands much harder, because the width stops looking like a choice somebody made.

\textbf{Spend the time in section 2}, which carries the misconception. Put both pictures up
before computing anything and ask which one is lying. Then work the \texttt{work} block live ---
both total 24 --- so nobody spends five minutes suspecting the data was tampered with. Do not
resolve the tension too fast: the sentence to land is that both are correct and they answer
different questions.

\textbf{Do not skip the counterweight} at the end of section~2 (blocks of one minute). Without
it students leave with ``narrower is better,'' which is wrong, and it is the version of this
error the AP exam actually tests.

\textbf{Guided Practice is where the pen changes hands.} \emph{The Bus Stop} is a dotplot, so
students meet the third display in a working context, and part (c) is the crux transferred: the
two far-out riders survive regrouping, but the empty stretch that made them look unusual does
not. If you only get through (a) and (c), that is the right trade.

\textbf{There is no independent practice set on this page, and that is deliberate --- the
homework is the individual practice.} Guided Practice is the last thing on the page; the period
ends with students starting the homework in front of you, which is where your second formative
read comes from. Item~2 is the one to watch: two histograms of the same 60 test scores, and
\emph{neither} is wrong. Sort as you circulate --- said neither and named the width (ready);
said neither with no reason (ready, needs the language); named one version wrong (the
misconception survived). If the third pile is more than a third of the room, \textbf{open
Lesson~1.4 by reteaching section~2}, and say so plainly.

\textbf{Debrief (8 min), spoken.} The two totals on the board first; it is the one move that
cannot be cut. \textbf{Do not let the debrief eat the last eight minutes} --- the homework start
is not optional padding.
\end{teachernote}

\begin{teachernote}[AP Practice]
\textbf{This page is not required and not scored} --- the graded practice for this lesson is the
homework behind it at the back of the packet. Work it as AP-format reps and go over it at the
start of Lesson~1.4. \textbf{MC item~2 is the one worth reading closely}: distractor (B) says
``discrete, because the recorded values are whole numbers,'' which is the exact reasoning half
the class used in the warm-up. If more than a handful pick it, spend three minutes on
measured-versus-counted before starting 1.4, because that confusion resurfaces in Unit~4 with
random variables. On the free response, part (d) is the highest-value part --- it forces students
to justify a display \emph{choice} from the variable type, which is Skill 2.B in the form the
exam actually asks it.
\end{teachernote}

\begin{teachernote}[Homework]
\textbf{Scored, 10 points (2 per item), due the first class after two study halls.} The eight minutes at the end of the
period are a \emph{start}, not a completion: put students on 1, 2, and~5 while you can still
catch a wrong start, and let 3 and~4 go home. Item~2 is the one to read first when scoring ---
``neither is wrong'' is the answer, and any student who names a wrong version has turned ``wide
hides things'' into a rule. Item~4 is the cheapest diagnostic in the set: a student who cannot
read 72, 74, 74, 79 off the stem is not going to survive Lesson~1.4. \textbf{When you score the
page, sort item~5 into three piles} --- chose (B) and justified it by the 200 values (ready);
chose (B) vaguely (ready, needs the language); chose (A) or (E) (the misconception survived). If
more than a third land in the last pile, reopen the two pictures in the Lesson~1.4 warm-up
debrief rather than letting it ride. Require item~3 to say what the horizontal axis is, not
merely that the bars touch.

\textbf{Paper is the default.} This page is what gets assigned unless you decide otherwise on a
given day. On the occasions you assign DeltaMath in its place, strike the score cell on the cover
and tell students this page is now optional practice. Do not assign both.
\end{teachernote}
\end{document}
